% State of the Art section. Paste into your own thesis, or \input it.
%
% Built from docs_thesis/find_research.md Stage 4.
% Verification status: shen2026geometry is full-text read (docs_thesis/papers/). The rest --
% sun2026angiographcad, sharp2026ssm, yaseliani2025gnn -- are [BIBLIOGRAPHY VERIFIED],
% abstract-level.

\section{The nearest competitors}
\label{sec:nearest-competitors}

Three lines of work sit closest to this thesis without doing what it does, and a fourth is
frequently mistaken for a fourth.

Sun et al.\ built AngioGraphCAD, a graph neural network over lesion geometry with masked attention
fusing multiple stenoses in the same patient, predicting future cardiovascular events at both the
lesion and the patient level from invasive coronary angiography \cite{sun2026angiographcad}.
Evaluated on two cohorts at lesion level and one at patient level, it outperformed clinical risk
measures at both, and the authors describe it as the first work to show geometric information
advancing event prediction from invasive angiography.
That claim has to be answered directly rather than ignored, and the differences that keep the two
projects distinct are concrete: invasive angiography rather than CCTA, the geometry of individual
stenoses rather than of the whole tree, and supervised event prediction rather than the population
description and unsupervised outlier detection objectives 8 and 9 ask for.
AngioGraphCAD strengthens the premise this thesis rests on, that geometry carries prognostic
information, while occupying a different problem.

Shen et al.\ published the review this field has been missing, written by several of the groups
behind the mechanism in Section~\ref{sec:shape-enough} \cite{shen2026geometry}.
It names two gaps directly relevant here: investigations relating coronary anatomy to hemodynamics
with a large population are limited, and most studies cover only the main bifurcation, when
vulnerable plaque sits mostly in the middle and distal vessels, so the whole tree needs investigating
rather than one junction.
Read carefully, its call for larger populations is a call to scale up patient-specific
computational-fluid-dynamics studies, not a call for outlier or extreme-value detection, and it
should not be read as the field already endorsing objective 9 --- only as the field naming objective
8's absence.

Sharp et al.\ reviewed statistical shape modeling across cardiovascular disease: landmark methods
through point distribution models, principal component analysis for dimensionality reduction, and
compactness, generalization and specificity as the standard metrics for evaluating a population shape
model \cite{sharp2026ssm}.
That vocabulary is the reference point if objective 9 takes a shape-model route to flagging outliers.
What the review does not contain matters as much as what it does: its subjects are cardiac chambers
and ventricles, and the coronary tree does not appear among them, so no part of that literature has
yet been applied to a coronary population.

Yaseliani et al.\ are recorded to head off a category error that recurs across this literature
\cite{yaseliani2025gnn}.
Their graph convolutional network predicts long-term CAD mortality at high accuracy, but its graph
has patients as nodes, linked by propensity-matched clinical features --- no artery is modeled at
all.
Most of what surfaces under "graph neural network for coronary artery disease" is this: a
patient-similarity graph, or a graph that labels centerline segments rather than one that models how
they connect into a tree.
That distinction has to be made explicitly, because it accounts for most of the apparent competition
this thesis has.
